\begin{longtable}{@{}l p{12.5cm}@{}}
  \caption{Content per-instance answer comparison ($\Delta \leq -75\%$) for the 8 most divergent items: same question, each model's answer and reasoning.}\label{tab:cmp_content_neg75} \\
  \toprule
  \endfirsthead
  \toprule
  \endhead
  \multicolumn{2}{@{}l}{\textbf{gpqa\_0080} \quad $\Delta = -100\%$} \\
  \textbf{Question} & Why does the hydroboration reaction between a conjugated diene and Ipc2BH form a single product, even at different temperatures? \newline A. The given reaction is stereospecific, and hence only one stereoisomer is formed. \newline B. The formation of the product is independent of the temperature at which the reaction takes place. \newline C. It is a concerted reaction, and no rearrangements are possible. \newline D. The reaction is syn-addition, which means both groups are added to the same face, leading to a single product. \\
  \textbf{Gold} & C \\
  \midrule
  \textbf{Kimi-K2.6} & answer C \\
  & The user is asking about the hydroboration of a conjugated diene with Ipc2BH (diisopinocampheylborane) and why it forms a single product even at different temperatures. \newline Let me think about the chemistry. Ipc2BH is a bulky, chiral hydroborating agent. It's used for asymmetric hydroboration. When reacting with a conjugated diene, the regioselectivity and stereoselectivity are important. \newline Wait, the question says "a conjugated diene" - but which one? Actually, Ipc2BH is known to react with certain dienes to give a single product. But the key here might be about the nature of the hydroboration with Ipc2BH. \newline Actually, I need to think about what makes the reaction give a single product regardless of temperature. \newline Let me analyze the options: \newline A. "The given reaction is stereospecific, and hence only one stereoisomer is formed." - Stereospecific means the stereochemistry of the reactant determines the stereochemistry of the product. But the question mentions different temperatures. If it's stereospecific, that doesn't necessarily explain why temperature doesn't affect the product distribution. \newline \textit{[\ldots\ truncated \ldots]} \newline Actually, let me reconsider option C's wording: "It is a concerted reaction, and no rearrangements are possible." This directly addresses why there's no temperature dependence: no intermediate, no rearrangement, so only the direct concerted product is formed. \newline Therefore, the answer is C. \\
  \textbf{Innospark-1T} & answer A \\
  & Let me analyze this question step by step. \newline The question asks: Why does the hydroboration reaction between a conjugated diene and Ipc2BH form a single product, even at different temperatures? \newline Let me consider each option: \newline - Option A: The given reaction is stereospecific, and hence only one stereoisomer is formed. \newline - Option B: The formation of the product is independent of the temperature at which the reaction takes place. \newline - Option C: It is a concerted reaction, and no rearrangements are possible. \newline \textit{[\ldots\ truncated \ldots]} \newline A. The given reaction is stereospecific, and hence only one stereoisomer is formed. \newline Based on this analysis, the answer is A. \\
  \specialrule{1.2pt}{8pt}{8pt}
  \multicolumn{2}{@{}l}{\textbf{cj\_eval\_0203} \quad $\Delta = -100\%$} \\
  \textbf{Question} & Draw a vertical map of China on a piece of paper that is 120 centimeters long and 120 centimeters wide. The most suitable scale is \newline A. 1 centimeter on the map equals an actual distance of 300 KM \newline B. 1:5000000 \newline C. 1/10000000 \newline D. one ten millionth \\
  \textbf{Gold} & B \\
  \midrule
  \textbf{Kimi-K2.6} & answer B \\
  & We need to answer a multiple-choice question about the most suitable scale for drawing a vertical map of China on a piece of paper 120 cm x 120 cm. \newline We need to determine the dimensions of China. China is roughly 5,500 km from north to south (vertical extent). Some sources say about 5,500 km north-south, and east-west about 5,200 km. Let's verify: The north-south extent of China is from about 53N (Mohe) to about 4N (Zengmu Ansha in Nansha Islands) but if we consider mainland China, from about 53N to about 18N (Hainan). Actually, the north-south distance of China is often cited as about 5,500 km. Some textbooks say 5,500 km. Let's assume the vertical map of China means the north-south extent. The paper is 120 cm long and 120 cm wide. So the map can be up to 120 cm in the vertical dimension. \newline We need to find a scale that fits China's north-south extent into 120 cm. \newline If north-south is 5,500 km = 550,000,000 cm. To fit into 120 cm, the scale would be 120 : 550,000,000 = 1 : (550,000,000/120) = 1 : 4,583,333. So roughly 1:4,600,000. \newline Let's check the options: \newline A. 1 cm on map equals 300 km actual. That means 1 cm : 300 km = 1 cm : 30,000,000 cm = 1:30,000,000. That would make the map length = 5500 km / 300 km per cm = 18.3 cm. That would fit easily, but maybe too small? The question asks "most suitable". If the paper is 120 cm, you might want to fill it. But maybe they want the largest possible scale that fits. Let's see other options. \newline \textit{[\ldots\ truncated \ldots]} \newline The question says "Answer with ONLY the letter (A, B, C, or D) of the correct answer. Do not include any explanation or additional text." \newline So I must output just "B". \\
  \textbf{Innospark-1T} & answer A \\
  & Let me analyze this question step by step. \newline The question asks: Draw a vertical map of China on a piece of paper that is 120 centimeters long and 120 centimeters wide. The most suitable scale is \newline Let me consider each option: \newline - Option A: 1 centimeter on the map equals an actual distance of 300 KM \newline - Option B: 1:5000000 \newline - Option C: 1/10000000 \newline \textit{[\ldots\ truncated \ldots]} \newline 120 cm scale 5,500 km. Converting 5,500 km to centimeters gives 5.5 10 cm, so the required scale is 120 / 5.5 10 1 : 4,583,333. Among the choices, only option A (1 cm = 300 km, i.e., 1 : 30,000,000) is small enough to fit the entire vertical span within 120 cm. Options B, C, and D all have larger scales and would require more than 120 cm to depict 5,500 km. Therefore, A is the correct answer. \newline Based on this analysis, I can determine the correct answer. \\
  \specialrule{1.2pt}{8pt}{8pt}
  \multicolumn{2}{@{}l}{\textbf{cj\_eval\_0489} \quad $\Delta = -100\%$} \\
  \textbf{Question} & Montesquieu occupies an important position in the French Enlightenment, and his theory of "separation of powers" has profoundly influenced the modern Western political democratic process and the entire history of humanity. Which of the following legal documents practices the principle of separation of powers and further improves it \newline A. "Napoleonic Code" \newline B. "Code of Hammurabi" \newline C. "U.S. Constitution of 1787" \newline D. "Provisional Constitution of the Republic of China" \\
  \textbf{Gold} & C \\
  \midrule
  \textbf{Kimi-K2.6} & answer C \\
  & The user wants me to answer a multiple-choice question about which legal document practices and further improves the principle of separation of powers (Montesquieu's theory). \newline Let me analyze each option: \newline A. "Napoleonic Code" (Civil Code of France) - This is primarily a civil code dealing with private law, property, family, etc. It doesn't primarily focus on separation of powers in government structure. \newline B. "Code of Hammurabi" - Ancient Babylonian law code, predates Enlightenment by millennia, has nothing to do with separation of powers. \newline C. "U.S. Constitution of 1787" - This is the key document. The U.S. Constitution explicitly implements separation of powers among legislative (Congress), executive (President), and judicial (Supreme Court) branches. It also adds checks and balances, which further improves/refines Montesquieu's theory by making the branches not just separate but also checking each other. The Founding Fathers were directly influenced by Montesquieu. \newline D. "Provisional Constitution of the Republic of China" - While it may have had some separation of powers elements, it was modeled after Western constitutions but is not the document most famous for practicing and improving upon Montesquieu's principle in the way the question implies. Also, historically it came much later (1912). \newline The question asks which document "practices the principle of separation of powers and further improves it." The U.S. Constitution of 1787 is the classic answer here because it established the three branches and the system of checks and balances, which is considered an improvement/refinement on Montesquieu's original theory. \newline Therefore, the correct answer is C. \\
  \textbf{Innospark-1T} & answer A \\
  & Let me analyze this question step by step. \newline The question asks: Montesquieu occupies an important position in the French Enlightenment, and his theory of "separation of powers" has profoundly influenced the modern Western political democratic process and the entire history of humanity. Which of the following legal documents practices the principle of separation of powers and further improves it \newline Let me consider each option: \newline - Option A: "Napoleonic Code" \newline - Option B: "Code of Hammurabi" \newline - Option C: "U.S. Constitution of 1787" \newline \textit{[\ldots\ truncated \ldots]} \newline The Constitution went further than Montesquieu's original conception by adding a system of checks and balances, where each branch has some control over the others, making it a more sophisticated implementation of the principle. D. "Provisional Constitution of the Republic of China" (1912): While this document did incorporate some separation of powers principles, it was heavily influenced by Western models including the U.S. However, it was short-lived and less influential in "further improving" the concept compared to the U.S. \newline Based on this analysis, the answer is C. \\
  \specialrule{1.2pt}{8pt}{8pt}
  \multicolumn{2}{@{}l}{\textbf{cj\_eval\_0036} \quad $\Delta = -100\%$} \\
  \textbf{Question} & Responding to external stimuli is one of the basic characteristics of living organisms. Which of the following options does not reflect this basic characteristic? \newline A. Antibodies are generated when the body is stimulated by external antigens \newline B. The sunflower head turns towards the sun under the guidance of sunlight \newline C. The pill bug curls up when it encounters external stimuli \newline D. Soybean seeds significantly swell after soaking \\
  \textbf{Gold} & D \\
  \midrule
  \textbf{Kimi-K2.6} & answer D \\
  & The user wants me to answer a multiple-choice question about biology. The question asks which option does NOT reflect the basic characteristic of living organisms responding to external stimuli. \newline Let me analyze each option: \newline A. Antibodies are generated when the body is stimulated by external antigens - This is an immune response to external stimuli (antigens). This is a response to external stimuli. \newline B. The sunflower head turns towards the sun under the guidance of sunlight - This is phototropism, a classic example of response to external stimuli (light). \newline C. The pill bug curls up when it encounters external stimuli - This is a defensive behavior in response to external stimuli (touch/threat). \newline D. Soybean seeds significantly swell after soaking - This is imbibition, a physical process where seeds absorb water and swell. This is not a physiological response to stimuli in the biological sense of irritability/sensitivity. It's a passive physical process of water absorption by osmosis/imbibition, not an active response characteristic of living organisms' irritability. While seeds are alive, the swelling itself is a physical/chemical process of water absorption, not a coordinated biological response to stimuli like the others. \newline \textit{[\ldots\ truncated \ldots]} \newline - D: Water absorption/imbibition = physical process, not a biological response to stimulus \newline Yes, D is the correct answer. \\
  \textbf{Innospark-1T} & answer A \\
  & Let me analyze this question step by step. \newline The question asks: Responding to external stimuli is one of the basic characteristics of living organisms. Which of the following options does not reflect this basic characteristic? \newline Let me consider each option: \newline - Option A: Antibodies are generated when the body is stimulated by external antigens \newline - Option B: The sunflower head turns towards the sun under the guidance of sunlight \newline - Option C: The pill bug curls up when it encounters external stimuli \newline \textit{[\ldots\ truncated \ldots]} \newline Option D involves no detection or response mechanism - it's purely a physical process of water absorption. \newline Based on this analysis, I can determine the correct answer. \\
  \specialrule{1.2pt}{8pt}{8pt}
  \multicolumn{2}{@{}l}{\textbf{cj\_eval\_0630} \quad $\Delta = -100\%$} \\
  \textbf{Question} & Use a ruler to measure the length of the physics textbook. Which of the following is a reason for measurement error \newline A. The ruler is not kept parallel to the long edge of the textbook \newline B. The value below the smallest scale is estimated inaccurately when reading \newline C. The ruler is not pressed against the textbook \newline D. The ruler is not aligned with the zero scale line \\
  \textbf{Gold} & B \\
  \midrule
  \textbf{Kimi-K2.6} & answer B \\
  & The user wants me to answer a multiple-choice question about measurement error when measuring the length of a physics textbook with a ruler. I need to choose the correct answer from A, B, C, or D. \newline Let me analyze each option: \newline A. The ruler is not kept parallel to the long edge of the textbook \newline - If the ruler is not parallel to the edge being measured, the measurement will be longer than the actual length (similar to the hypotenuse being longer than the side). This is indeed a source of measurement error. \newline B. The value below the smallest scale is estimated inaccurately when reading \newline - This refers to estimation error when reading between the smallest divisions (the doubtful digit). This is a random error that occurs in every measurement and is indeed a reason for measurement error. \newline \textit{[\ldots\ truncated \ldots]} \newline Yes, this makes sense. The question asks for "measurement error" , and in the context of Chinese physics, estimation inaccuracy is the classic example of unavoidable measurement error, while the others are operational mistakes that should be avoided. \newline So the correct answer is B. \\
  \textbf{Innospark-1T} & answer A \\
  & Let me analyze this question step by step. \newline The question asks: Use a ruler to measure the length of the physics textbook. Which of the following is a reason for measurement error \newline Let me consider each option: \newline - Option A: The ruler is not kept parallel to the long edge of the textbook \newline - Option B: The value below the smallest scale is estimated inaccurately when reading \newline - Option C: The ruler is not pressed against the textbook \newline \textit{[\ldots\ truncated \ldots]} \newline Option A represents a clear geometric measurement error where the ruler's misalignment causes the measured path to be longer than the actual object dimension. This is a fundamental concept in measurement science - the ruler must be parallel to the dimension being measured to obtain an accurate length. \newline Based on this analysis, the answer is A. \\
  \specialrule{1.2pt}{8pt}{8pt}
  \multicolumn{2}{@{}l}{\textbf{cj\_eval\_0773} \quad $\Delta = -100\%$} \\
  \textbf{Question} & Uncooked green shrimp turn red after being cooked. Some students believe that this red substance may be similar to an acid-base indicator, changing color when it encounters an acid or a base. According to these students' views, this falls under in scientific inquiry. \newline A. Observation \newline B. Experiment \newline C. Hypothesis \newline D. Conclusion \\
  \textbf{Gold} & C \\
  \midrule
  \textbf{Kimi-K2.6} & answer C \\
  & The user wants me to answer a multiple-choice question about scientific inquiry. \newline The question describes: Uncooked green shrimp turn red after being cooked. Some students believe that this red substance may be similar to an acid-base indicator, changing color when it encounters an acid or a base. According to these students' views, this falls under in scientific inquiry. \newline Options: \newline A. Observation \newline B. Experiment \newline C. Hypothesis \newline \textit{[\ldots\ truncated \ldots]} \newline So the students' belief that the red substance may be similar to an acid-base indicator is a hypothesis. \newline Answer: C \\
  \textbf{Innospark-1T} & answer A \\
  & Let me analyze this question step by step. \newline The question asks: Uncooked green shrimp turn red after being cooked. Some students believe that this red substance may be similar to an acid-base indicator, changing color when it encounters an acid or a base. According to these students' views, this falls under in scientific inquiry. \newline Let me consider each option: \newline - Option A: Observation \newline - Option B: Experiment \newline - Option C: Hypothesis \newline \textit{[\ldots\ truncated \ldots]} \newline This is the essence of a hypothesis in scientific inquiry. \newline Based on this analysis, the answer is C. \\
  \specialrule{1.2pt}{8pt}{8pt}
  \multicolumn{2}{@{}l}{\textbf{cj\_eval\_0553} \quad $\Delta = -100\%$} \\
  \textbf{Question} & Which pair of options represents opposite meanings? \newline A. Winning 3 games and losing 2 games \newline B. Offense and defense \newline C. Temperature rising by 3 and temperature dropping to 2 \newline D. Increasing income by 30,000 yuan and decreasing expenditure by 20,000 yuan \\
  \textbf{Gold} & A \\
  \midrule
  \textbf{Kimi-K2.6} & answer A \\
  & A \\
  \textbf{Innospark-1T} & answer B \\
  & Let me analyze each option to determine which pair represents opposite meanings: \newline **Option A: Winning 3 games and losing 2 games** \newline - These are opposite concepts (winning vs. losing) \newline - However, the quantities (3 games vs. 2 games) are different \newline - This is not a clean pair of opposites due to the mismatched numbers \newline **Option B: Offense and defense** \newline \textit{[\ldots\ truncated \ldots]} \newline **Conclusion:** \newline Option B (Offense and defense) represents the clearest pair of opposite meanings. In competitive contexts, offense and defense are fundamentally opposing concepts - one focuses on attacking/scoring while the other focuses on protecting/preventing scores. \\
  \specialrule{1.2pt}{8pt}{8pt}
  \multicolumn{2}{@{}l}{\textbf{gpqa\_0045} \quad $\Delta = -100\%$} \\
  \textbf{Question} & Consider the Y-component of the intrinsic angular momentum operator, A of a muon be represented by a 22 matrix Ay satisfying the eigenvalue equation Ay = a. Here, a is the eigenvalue, is the eigenfunction. The matrix operator has the form Ay = cS; where the constant c=h/4 and S being a 22 matrix. The first row of the matrix S is (0 -i) and the second row is (i 0). You are asked to calculate the eigenvalue and eigenvectors of the operator Ay. During the calculation, which statement below will you consider to be correct? \newline A. The imaginary part of the eigenvalue of Ay are +1/2 or --1/2, and the real part of that are +1 or --1. \newline B. The imaginary part of the eigenvalue of Ay are +2h or --2h, and the real part of that are +h/4 or --h/4. \newline C. The eigenfunction of the operator Ay can also be an eigenfunction of A\textasciicircum{}2, but not of the Z-component, Az. \newline D. The eigenfunctions of the operator Ay are the basis functions of the matrix operator Ay given above. \\
  \textbf{Gold} & C \\
  \midrule
  \textbf{Kimi-K2.6} & answer C \\
  & We need to parse the problem carefully. The question is about the Y-component of the intrinsic angular momentum operator of a muon, represented by a 2x2 matrix Ay satisfying Ay = a . Ay = c S, where c = h/(4) and S is a 2x2 matrix with first row (0, -i) and second row (i, 0). So S is the Pauli matrix \_y (the standard Pauli matrix for y-component). Indeed \_y = [[0, -i], [i, 0]]. Then Ay = (h/(4)) \_y. \newline We need to calculate eigenvalues and eigenvectors of Ay. Then choose which statement is correct. \newline First, eigenvalues of \_y are +1 and -1. Because \_y is Hermitian, trace 0, determinant 1, eigenvalues 1. So eigenvalues of Ay = c \_y are c*(1) = h/(4). So eigenvalues are real numbers: h/(4). There is no imaginary part. So statements A and B are wrong because they talk about imaginary parts. Let's check them: \newline A. "The imaginary part of the eigenvalue of Ay are +1/2 or --1/2, and the real part of that are +1 or --1." This is nonsense; eigenvalues are real. \newline B. "The imaginary part of the eigenvalue of Ay are +2h or --2h, and the real part of that are +h/4 or --h/4." Also nonsense. \newline C. "The eigenfunction of the operator Ay can also be an eigenfunction of A\textasciicircum{}2, but not of the Z-component, Az." \newline \textit{[\ldots\ truncated \ldots]} \newline Therefore, final answer: C. \newline The user wants ONLY the letter. So output "C". \\
  \textbf{Innospark-1T} & answer D \\
  & Let me analyze this question step by step. \newline The question asks: Consider the Y-component of the intrinsic angular momentum operator, A of a muon be represented by a 22 matrix Ay satisfying the eigenvalue equation Ay = a. Here, a is the eigenvalue, is the eigenfunction. The matrix operator has the form Ay = cS; where the constant c=h/4 and S being a 22 \newline Let me consider each option: \newline - Option A: The imaginary part of the eigenvalue of Ay are +1/2 or --1/2, and the real part of that are +1 or --1. \newline - Option B: The imaginary part of the eigenvalue of Ay are +2h or --2h, and the real part of that are +h/4 or --h/4. \newline - Option C: The eigenfunction of the operator Ay can also be an eigenfunction of A\textasciicircum{}2, but not of the Z-component, Az. \newline \textit{[\ldots\ truncated \ldots]} \newline Final Answer: D \newline Based on this analysis, the answer is D. \\
  \specialrule{1.2pt}{8pt}{8pt}
  \bottomrule
\end{longtable}